% ===========================================================================
% CAPÍTULO 4 — RESULTADOS: validación de la placa de comunicación serie
% ===========================================================================

\section{Validación de la placa de comunicación serie} \label{sec:validacion_propia}

\subsection{Prueba previa: 14 transceptores simultáneos}

Para verificar el correcto funcionamiento del driver con un número elevado de instancias, se
cargó un \textit{bitstream} con 14 transceptores instanciados (canales 0 a 13), conectando los
canales adicionales a los pines externos del conector J3 de la ZCU102. El sistema arrancó y operó
correctamente con todos los canales, confirmando que la arquitectura de ISR maestra con tareas
\textit{worker} individuales escala sin problemas hasta el número máximo de instancias soportado.
La prueba con los 14 transceptores sobre la placa de comunicación serie ya
fabricada se aborda en la sección~\ref{sec:resultados_pcb}.

\subsection{Comparativa de las arquitecturas de transporte PS--PL} \label{sec:resultados_benchmark}

Esta sección recoge la medida de las tres arquitecturas descritas en la
sección~\ref{sec:transporte_desarrollo}. Las tres se ejecutaron sobre la misma ZCU102, con el
mismo transceptor serie de catorce canales, el mismo programa de pruebas
(\texttt{bench\_main.c}) y los buses cerrados por el módulo de \textit{loopback} de la PL
descrito en la sección~\ref{sec:bench_loopback}. Lo único que difiere entre las tres corridas es
el transporte y su driver, que es precisamente lo que se quiere comparar.

\paragraph{Una salvedad previa.} En la variante DMA14 \textbf{no se consiguió hacer funcionar el
módulo de \textit{loopback}}: su autodiagnóstico detectó que el lazo de recepción estaba abierto
y, en consecuencia, no llegó ningún byte. Como se explica en la
sección~\ref{sec:transporte_desarrollo}, se decidió no invertir más tiempo en depurarlo, porque
el coste en lógica y el límite de líneas de interrupción ya descartaban esta arquitectura como
solución final. Todo lo que depende de recibir ---latencia, \textit{throughput}, multipunto,
desbordamiento, recuperación--- sale a cero en esa variante, y ese cero \textbf{no mide el
transporte}: mide la ausencia de lazo. Se hace constar con un guion en las tablas. Lo que sí es
plenamente válido de esa corrida, y lo que sostiene el argumento sobre esta variante, es el coste
de hardware, la huella de memoria y el coste en interrupciones de la transmisión.

% ---------------------------------------------------------------------------
\subsubsection{Coste de hardware}
% ---------------------------------------------------------------------------

Las tres arquitecturas están implementadas sobre el mismo dispositivo (XCZU9EG, 274\,080 LUT,
548\,160 registros y 912 bloques de RAM), de modo que las cifras salen directamente de los
informes de utilización de Vivado y son comparables sin normalización alguna.

\begin{table}[H]
    \centering
    \begin{tabular}{lrrrrrrr}
        \hline
        \textbf{Variante} & \textbf{LUT} & \textbf{\%} & \textbf{Registros} & \textbf{BRAM} & \textbf{Potencia} & \textbf{WNS} & \textbf{LUT/canal} \\
        \hline
        GPIO   & 6\,434  & 2,35\,\%  & 7\,726  & 7  & 3,52 W & 4,69 ns & 460 \\
        DMA14  & 40\,217 & 14,67\,\% & 53\,943 & 35 & 3,86 W & 4,05 ns & 2\,873 \\
        MCDMA  & 14\,191 & 5,18\,\%  & 16\,438 & 18 & 3,62 W & 5,14 ns & 1\,014 \\
        \hline
    \end{tabular}
    \caption{Ocupación de la lógica programable de las tres variantes sobre el XCZU9EG.}
    \label{tab:bench_hw}
\end{table}

La figura~\ref{fig:bench_hw} representa las tres magnitudes normalizadas como porcentaje de los
recursos disponibles en el dispositivo, que es lo que permite dibujarlas sobre un mismo eje y
compararlas de un vistazo.

\begin{figure}[H]
    \centering
    \begin{tikzpicture}
        \begin{axis}[
            benchbar,
            height=6.2cm,
            ylabel={\% de los recursos del XCZU9EG},
            ymax=17,
            symbolic x coords={LUT, Registros, BRAM},
            nodes near coords={\pgfmathprintnumber[fixed, fixed zerofill, precision=2]{\pgfplotspointmeta}\,\%},
            legend style={at={(0.5,1.06)}, anchor=south, legend columns=3},
        ]
        \addplot[fill=cGPIO,  draw=none] coordinates {(LUT,2.35) (Registros,1.41) (BRAM,0.77)};
        \addplot[fill=cDMA14, draw=none] coordinates {(LUT,14.67) (Registros,9.84) (BRAM,3.84)};
        \addplot[fill=cMCDMA, draw=none] coordinates {(LUT,5.18) (Registros,3.00) (BRAM,1.97)};
        \legend{GPIO, DMA14, MCDMA}
        \end{axis}
    \end{tikzpicture}
    \caption{Ocupación de la FPGA por variante, en porcentaje de los recursos del dispositivo.
    Menos es mejor. El banco de catorce DMA triplica largamente el consumo de la solución
    multicanal para el mismo número de canales.}
    \label{fig:bench_hw}
\end{figure}

La cifra que resume la comparación es la última columna de la tabla. El banco de catorce AXI DMA cuesta
\textbf{2,8 veces más lógica que el MCDMA} para exactamente el mismo número de canales, porque
replica catorce motores completos en lugar de compartir uno multicanal. La variante GPIO es, con
diferencia, la más pequeña ---y así debe ser: apenas contiene los transceptores y un puñado de
registros---, pero ese ahorro de lógica no es gratuito: lo paga la CPU, como se ve más abajo.

El margen de \textit{timing} (WNS) es positivo y holgado en las tres, de modo que ninguna cierra
al límite y las diferencias de rendimiento no son atribuibles a la implementación física.

% ---------------------------------------------------------------------------
\subsubsection{Coste de software}
% ---------------------------------------------------------------------------

\begin{table}[H]
    \centering
    \begin{tabular}{lrrrrr}
        \hline
        \textbf{Variante} & \textbf{Código} & \textbf{RAM del driver} & \textbf{Líneas IRQ} & \textbf{Tareas} & \textbf{Semáforos} \\
        \hline
        GPIO  & 3\,099 B & 114\,688 B & 1 & 14 & 28 \\
        DMA14 & 3\,036 B & 115\,584 B & 2 & 1  & 15 \\
        MCDMA & 3\,884 B & 122\,434 B & 3 & 1  & 17 \\
        \hline
    \end{tabular}
    \caption{Huella de software del driver de cada variante (tamaño de \texttt{.text} y recursos RTEMS reservados en ejecución).}
    \label{tab:bench_sw}
\end{table}

La memoria total reservada es prácticamente la misma en las tres ---en torno a 112--120~KB,
dominada por los \textit{buffers} circulares de 4~KB por canal---, aunque su origen difiere: la
variante GPIO los pide con \texttt{malloc} y las de DMA los reservan estáticamente, porque el
motor necesita direcciones fijas y alineadas.

La diferencia significativa está en las \textbf{primitivas de RTEMS}. La variante GPIO necesita
\textbf{una tarea de recepción por canal} ---catorce tareas y veintiocho semáforos--- porque cada
canal drena su propio FIFO hardware de forma independiente. Las dos variantes con DMA se apañan
con \textbf{una sola tarea de despacho}, ya que el motor entrega paquetes completos y la ISR solo
tiene que señalizar quién ha terminado. Trece tareas menos es trece cambios de contexto menos, y
trece pilas menos.

% ---------------------------------------------------------------------------
\subsubsection{Coste en interrupciones}
% ---------------------------------------------------------------------------

Esta es la métrica que explica todas las demás. Cuenta cuántas veces se expropia la CPU por cada
kilobyte movido, acumulado sobre la sesión completa de pruebas.

\begin{table}[H]
    \centering
    \begin{tabular}{lrrrrr}
        \hline
        \textbf{Variante} & \textbf{IRQ/KB} & \textbf{IRQ TX} & \textbf{IRQ RX} & \textbf{Bytes TX} & \textbf{Bytes RX} \\
        \hline
        GPIO  & 1\,250 & 22\,161 & 129\,382 & 22\,161 & 101\,890 \\
        DMA14 & 3\,$^{*}$ & 85 & --- & 23\,132 & --- \\
        MCDMA & \textbf{3} & 194 & 192 & 23\,132 & 101\,472 \\
        \hline
        \multicolumn{6}{l}{\footnotesize $^{*}$ Solo transmisión: sin lazo de recepción, la cifra no es comparable.} \\
    \end{tabular}
    \caption{Interrupciones por kilobyte transferido.}
    \label{tab:bench_irq}
\end{table}

\begin{figure}[H]
    \centering
    \begin{tikzpicture}
        \begin{axis}[
            benchbar,
            height=5.6cm,
            width=0.76\textwidth,
            ylabel={Interrupciones por KB transferido},
            bar shift=0pt,
            ymax=1500,
            bar width=30pt,
            symbolic x coords={GPIO, DMA14, MCDMA},
            xtick={GPIO, DMA14, MCDMA},
            nodes near coords={\pgfmathprintnumber[fixed]{\pgfplotspointmeta}},
        ]
        \addplot[fill=cGPIO, draw=none] coordinates {(GPIO,1250)};
        \addplot[fill=cDMA14, draw=none] coordinates {(DMA14,3)};
        \addplot[fill=cMCDMA, draw=none] coordinates {(MCDMA,3)};
        \end{axis}
    \end{tikzpicture}
    \caption{Interrupciones por kilobyte transferido. Menos es mejor: cada interrupción es una
    expropiación de la CPU. Las barras de DMA14 y MCDMA son invisibles a esta escala, y eso es
    precisamente el resultado. La cifra de DMA14 recoge solo la transmisión.}
    \label{fig:bench_irq}
\end{figure}

Los números hablan por sí solos. En la variante GPIO, el número de interrupciones de transmisión
\textbf{coincide exactamente con el número de bytes transmitidos} (22\,161 y 22\,161): es la
confirmación empírica de que el procesador toca cada byte. En recepción ocurre lo mismo, con un
sobrecoste añadido por los eventos de estado. El resultado son \textbf{1250 interrupciones por
kilobyte}.

El MCDMA mueve la misma cantidad de datos con \textbf{194 interrupciones de transmisión y 192 de
recepción}: tres por kilobyte, unas \textbf{cuatrocientas veces menos}. La CPU pasa de estar
atendiendo el bus a estar disponible para la aplicación.

% ---------------------------------------------------------------------------
\subsubsection{Latencia}
% ---------------------------------------------------------------------------

Tiempo desde justo antes de la llamada de envío hasta que el receptor tiene el frame de 11 bytes
completo, sobre 50 muestras.

\begin{table}[H]
    \centering
    \begin{tabular}{lrrrrr}
        \hline
        \textbf{Variante} & \textbf{Media ($\mu$s)} & \textbf{p50} & \textbf{p95} & \textbf{Máx.} & \textbf{Timeouts} \\
        \hline
        GPIO  & 4\,606 & 4\,608 & 4\,611 & 4\,611 & 0 \\
        DMA14 & --- & --- & --- & --- & 50 \\
        MCDMA & \textbf{1\,530} & 1\,529 & 1\,535 & 1\,535 & 0 \\
        \hline
    \end{tabular}
    \caption{Latencia de entrega de un frame de 11 bytes, del envío en el maestro a la
    recepción completa en el esclavo (115\,200 8N1).}
    \label{tab:bench_latencia}
\end{table}

\begin{figure}[H]
    \centering
    \begin{tikzpicture}
        \begin{axis}[
            benchbar,
            height=5.6cm,
            width=0.7\textwidth,
            ylabel={Latencia media ($\mu$s)},
            bar shift=0pt,
            ymax=5600,
            scaled y ticks=false,
            ytick={0,1000,2000,3000,4000,5000},
            bar width=34pt,
            symbolic x coords={GPIO, MCDMA},
            xtick={GPIO, MCDMA},
            nodes near coords={\pgfmathprintnumber[fixed]{\pgfplotspointmeta}},
        ]
        \addplot[fill=cGPIO, draw=none] coordinates {(GPIO,4606)};
        \addplot[fill=cMCDMA, draw=none] coordinates {(MCDMA,1530)};
        \end{axis}
    \end{tikzpicture}
    \caption{Latencia media de entrega de un frame de 11 bytes. Menos es mejor. Se excluye
    DMA14 por carecer de lazo de recepción.}
    \label{fig:bench_latencia}
\end{figure}

El MCDMA es \textbf{tres veces más rápido} que el GPIO, y con una dispersión muy baja: la
diferencia entre la mediana y el máximo es de 6~$\mu$s sobre 50 muestras, lo que indica un
comportamiento notablemente determinista ---una propiedad especialmente valiosa en un sistema de
tiempo real embarcado.

% ---------------------------------------------------------------------------
\subsubsection{Throughput}
% ---------------------------------------------------------------------------

A 115\,200 baudios en formato 8N1, cada byte ocupa 10 bits de línea, de modo que el techo físico
de un canal es de \textbf{11\,520 B/s}. Esa es la referencia contra la que hay que leer la
primera columna.

\begin{table}[H]
    \centering
    \begin{tabular}{lrrrc}
        \hline
        \textbf{Variante} & \textbf{1 canal} & \textbf{3 maestros} & \textbf{6 receptores} & \textbf{Reps. sin pérdidas} \\
        \hline
        GPIO  & 3\,471 (30\,\%) & 10\,411 & 20\,838 & 5 / 5 \\
        DMA14 & --- & --- & --- & --- \\
        MCDMA & \textbf{11\,461 (99\,\%)} & \textbf{34\,303} & \textbf{68\,504} & 5 / 5 \\
        \hline
    \end{tabular}
    \caption{\textit{Throughput} en B/s. El porcentaje es sobre el techo físico de la línea (11\,520 B/s).}
    \label{tab:bench_throughput}
\end{table}

\begin{figure}[H]
    \centering
    \begin{tikzpicture}
        \begin{axis}[
            benchbar,
            height=6cm,
            width=0.7\textwidth,
            ylabel={\textit{Throughput} de un canal (B/s)},
            bar shift=0pt,
            ymax=14000,
            scaled y ticks=false,
            ytick={0,2000,4000,6000,8000,10000,12000},
            % halo blanco tras la cifra: la línea de techo pasa justo por detrás
            every node near coord/.append style={anchor=south, yshift=1pt,
                fill=white, inner sep=1.2pt},
            bar width=34pt,
            symbolic x coords={GPIO, MCDMA},
            xtick={GPIO, MCDMA},
            nodes near coords={\pgfmathprintnumber[fixed]{\pgfplotspointmeta}},
        ]
        % Techo físico de la línea a 115200 8N1 (debajo de las barras y sus etiquetas)
        \draw[cRef, dashed, thick] (axis cs:GPIO,11520) -- (axis cs:MCDMA,11520);
        \node[anchor=south west, font=\footnotesize, text=black!65]
            at (axis cs:GPIO,11700) {techo de línea: 11\,520 B/s};
        \addplot[fill=cGPIO, draw=none] coordinates {(GPIO,3471)};
        \addplot[fill=cMCDMA, draw=none] coordinates {(MCDMA,11461)};
        \end{axis}
    \end{tikzpicture}
    \caption{\textit{Throughput} de un canal frente al techo físico de la línea a 115\,200 8N1.
    Más es mejor. El MCDMA lo alcanza al 99\,\%; la variante GPIO se queda en el 30\,\%.}
    \label{fig:bench_throughput}
\end{figure}

El MCDMA alcanza el \textbf{99\,\% del techo físico de la línea}: el transporte deja de ser el
cuello de botella y el único límite pasa a ser la propia línea serie, que es exactamente el
objetivo de diseño. La variante GPIO se queda en el 30\,\%: dos de cada tres bytes que la línea
podría llevar se pierden atendiendo interrupciones.

Con carga concurrente la diferencia se amplía en lugar de reducirse: con seis receptores
simultáneos en el bus multipunto, el MCDMA sostiene \textbf{68\,504~B/s} frente a los 20\,838 del
GPIO, y lo hace sin perder un solo byte.

% ---------------------------------------------------------------------------
\subsubsection{Barrido de velocidad}
% ---------------------------------------------------------------------------

La prueba más reveladora de todas. Se reprograman los catorce transceptores a cada velocidad y se
mide un bloque de 1~KB en cada punto.

\begin{table}[H]
    \centering
    \begin{tabular}{lrrrrrrr}
        \hline
        \textbf{Variante} & \textbf{115200} & \textbf{230400} & \textbf{460800} & \textbf{921600} & \textbf{1M} & \textbf{2M} & \textbf{4M} \\
        \hline
        GPIO  & 3\,471 & 4\,087 & 4\,485 & 4\,714 & 4\,733 & 4\,848 & 4\,907 \\
        MCDMA & 11\,458 & 22\,789 & 45\,104 & 88\,298 & 96\,186 & 183\,053 & \textbf{344\,781} \\
        \hline
    \end{tabular}
    \caption{\textit{Throughput} (B/s) frente a la velocidad de línea (baudios). La variante DMA14 se omite por carecer de lazo de recepción.}
    \label{tab:bench_baudios}
\end{table}

\begin{figure}[H]
    \centering
    \begin{tikzpicture}
        \begin{axis}[
            bench,
            height=7.4cm,
            xmode=log, ymode=log,
            log basis x=10, log basis y=10,
            xlabel={Velocidad de línea (baudios)},
            ylabel={\textit{Throughput} (B/s)},
            xmin=90000, xmax=5500000,
            ymin=2000,  ymax=600000,
            xtick={115200, 230400, 460800, 921600, 2000000, 4000000},
            xticklabels={115k, 230k, 461k, 922k, 2M, 4M},
            xmajorgrids=true,
            legend style={at={(0.02,0.98)}, anchor=north west},
        ]
        \addplot[color=cMCDMA, line width=1.4pt, mark=*, mark size=2.4pt,
                 mark options={fill=cMCDMA, draw=white, line width=0.6pt}]
            coordinates {
                (115200,11458) (230400,22789) (460800,45104) (921600,88298)
                (1000000,96186) (2000000,183053) (4000000,344781)};
        \addlegendentry{MCDMA}

        \addplot[color=cGPIO, line width=1.4pt, mark=square*, mark size=2.2pt,
                 mark options={fill=cGPIO, draw=white, line width=0.6pt}]
            coordinates {
                (115200,3471) (230400,4087) (460800,4485) (921600,4714)
                (1000000,4733) (2000000,4848) (4000000,4907)};
        \addlegendentry{GPIO}

        % Anotaciones colocadas en hueco libre, con halo blanco para no pisar las curvas.
        \node[anchor=west, font=\footnotesize, text=cMCDMA!80!black,
              fill=white, inner sep=1.5pt]
            at (axis cs:1250000,32000) {escala con la línea};
        \node[anchor=north west, font=\footnotesize, text=cGPIO!85!black,
              fill=white, inner sep=1.5pt]
            at (axis cs:250000,3750) {saturado: el cuello de botella es la CPU};
        \end{axis}
    \end{tikzpicture}
    \caption{\textit{Throughput} frente a la velocidad de línea, en escala logarítmica en ambos
    ejes. El MCDMA crece de forma proporcional a la velocidad del enlace; la variante GPIO es una
    recta plana, porque su límite no está en la línea sino en el coste de una interrupción por
    byte. En el punto de 4~Mbaudios la diferencia es de 70 veces.}
    \label{fig:bench_baudios}
\end{figure}

Aquí se ve, sin necesidad de más argumentos, la naturaleza del límite de cada arquitectura:

\begin{itemize}
    \item El \textbf{GPIO se satura}. Multiplicar la velocidad de línea por 35 (de 115\,200 a
    4~Mbaudios) mejora el \textit{throughput} en un mísero 41\,\%, de 3471 a 4907~B/s. La curva es
    plana porque \textbf{el cuello de botella no es la línea, es la CPU}: el coste de una
    interrupción por byte no depende de lo rápido que vaya el bus. Por rápido que se haga el
    hardware serie, el sistema no aprovechará esa velocidad.

    \item El \textbf{MCDMA escala linealmente}. El \textit{throughput} crece de forma
    prácticamente proporcional a la velocidad de línea, hasta \textbf{344\,781~B/s} a
    4~Mbaudios: \textbf{70 veces} lo que consigue el GPIO en el mismo punto. La curva sigue a la
    línea porque el transporte ya no interviene.
\end{itemize}

En los siete puntos, y en ambas variantes, el número de errores de transporte es \textbf{cero}: la
saturación del GPIO no produce corrupción, solo pérdida de rendimiento.

% ---------------------------------------------------------------------------
\subsubsection{Robustez}
% ---------------------------------------------------------------------------

Dos pruebas adicionales verifican que la ganancia de rendimiento no se paga con fragilidad: el
envío de tres frames corruptos seguidos de uno válido (T7) y el envío de 8~KB sin que la
aplicación vacíe un \textit{ring} de 4~KB (T5), diseñado explícitamente para desbordar.

\begin{table}[H]
    \centering
    \begin{tabular}{lccrr}
        \hline
        \textbf{Variante} & \textbf{Frames corruptos} & \textbf{Se recupera} & \textbf{Bytes aceptados} & \textbf{Bytes descartados} \\
        \hline
        GPIO  & 3 / 3 & sí & 24\,576 & 16\,246 \\
        DMA14 & 3 / 3 & --- & --- & --- \\
        MCDMA & 3 / 3 & sí & 24\,576 & 24\,576 \\
        \hline
    \end{tabular}
    \caption{Rechazo de frames corruptos (T7) y comportamiento en desbordamiento del \textit{ring} de software (T5).}
    \label{tab:bench_robustez}
\end{table}

Ambas variantes rechazan los tres frames corruptos y siguen aceptando el frame válido posterior:
el error no deja el driver en un estado inconsistente.

La columna de bytes descartados merece una aclaración, porque se presta a leerse al revés. \textbf{No
son errores de transporte} ---esos son cero en las tres variantes---: cuentan un byte cada vez que
un dato llega correctamente pero \textbf{no cabe en el \textit{ring} de software porque nadie lo ha
leído}. La mayor parte procede de la prueba T5, cuyo propósito es justamente ese, y el resto de los
canales que nadie lee en las pruebas de un solo maestro: al ser un bus multipunto, los seis
receptores reciben cada bloque.

Por eso el número del MCDMA es \textit{mayor} que el del GPIO, y eso es una buena noticia y no una
mala: en la variante GPIO el propio cuello de botella de interrupciones frena la entrega ---el
\textit{ring} se llena más despacio y desborda menos---, mientras que el MCDMA entrega a la
velocidad de la línea todo lo que ésta produce y, en consecuencia, sobra más en los
\textit{buffers} que la aplicación de prueba deliberadamente no vacía. \textbf{Un contador de
descartes alto es aquí la firma de un transporte sano, no de uno defectuoso.} Conviene recordar
que estas cifras corresponden a la campaña final, con las treinta y seis comprobaciones en verde:
las corridas anteriores a las correcciones de contrapresión y de tamaño de paquete descritas en la
sección~\ref{sec:transporte_depuracion} entregaban menos datos a los \textit{rings} y no son
comparables.

% ---------------------------------------------------------------------------
\subsubsection{Discusión: qué arquitectura y por qué}
% ---------------------------------------------------------------------------

\begin{table}[H]
    \centering
    \begin{tabular}{p{3.2cm} p{3.4cm} p{3.4cm} p{3.4cm}}
        \hline
         & \textbf{GPIO} & \textbf{DMA14} & \textbf{MCDMA} \\
        \hline
        Lógica (LUT) & \textbf{6\,434} & 40\,217 & 14\,191 \\
        IRQ/KB & 1\,250 & 3 (solo TX) & \textbf{3} \\
        Throughput 1 canal & 3\,471 B/s (30\,\%) & --- & \textbf{11\,461 B/s (99\,\%)} \\
        Latencia & 4\,606 $\mu$s & --- & \textbf{1\,530 $\mu$s} \\
        Escala con el baudrate & no & --- & \textbf{sí} \\
        Líneas IRQ necesarias & 1 & 28 (reducidas a 2) & 3 \\
        Complejidad de diseño & \textbf{baja} & media & alta \\
        \hline
    \end{tabular}
    \caption{Síntesis de la comparativa. En negrita, el mejor valor de cada fila.}
    \label{tab:bench_resumen}
\end{table}

La conclusión que sostienen los datos es que \textbf{el MCDMA con puente VHDL es la arquitectura
adecuada para este sistema}, y que su ventaja no procede de ser «más rápido» en un sentido vago,
sino de una causa concreta y medible: \textbf{mueve paquetes en lugar de bytes}. Las 1250
interrupciones por kilobyte de la variante GPIO se convierten en 3, y todo lo demás ---la
latencia tres veces menor, el 99\,\% del techo de línea, la escalabilidad con la velocidad--- es
consecuencia de eso.

La variante DMA14 es instructiva precisamente por lo que demuestra que \textbf{no} hay que hacer.
Fue el primer intento de mejorar la arquitectura frente al AXI GPIO, y resuelve el problema
correcto ---quitar a la CPU de en medio--- pero por el camino equivocado: replicar catorce
motores cuesta 2,8 veces más lógica y, sobre todo, exige veintiocho líneas de interrupción cuando
el silicio solo ofrece ocho. La arquitectura no escala, y el límite no es del diseño sino del
dispositivo. A ello se suma un argumento de futuro: un 14,7\,\% del dispositivo es una hipoteca
incompatible con la triplicación TMR que exigiría endurecer el diseño frente a la radiación,
mientras que el 5,2\,\% del MCDMA deja ese margen abierto. Es el conjunto de razones que
justificó dejar de invertir en esta variante cuando su lazo de recepción se resistió, y el que
sostiene el esfuerzo de diseñar el puente VHDL de la variante C: no era una optimización, era la
única vía.

En cuanto al GPIO, sigue siendo una opción perfectamente razonable \textbf{para pocos canales a
baja velocidad}: es cuatro veces más pequeño en lógica, no necesita mapeo de MMU ni gestión de
coherencia de caché, y su driver es sensiblemente más simple. Lo que los datos demuestran es que
\textbf{no sirve para catorce canales}, y que su limitación no se cura subiendo la velocidad de la
línea, porque el cuello de botella está en el otro extremo.
